% META: {"id": "1NSI-AGT-EX-003", "chapitre": "1NSI-ALGO-PARCOURS-TRIS", "type_objet": "exercice", "capacites": ["P-ALGO-02A", "P-ALGO-02C"], "parcours": 2, "competences": ["calculer", "raisonner"], "duree_min": 15, "mode_creation": "ex_nihilo", "sources_inspiration": [], "fichier_tex": "chapitres/1NSI-ALGO-PARCOURS-TRIS/exercices/1NSI-AGT-EX-003.tex", "corrige_tex": "chapitres/1NSI-ALGO-PARCOURS-TRIS/corriges/1NSI-AGT-CO-003.tex", "status": "generated"}
% BEGIN-VERIFY
% def tri_insertion(tableau):
%     tableau = tableau[:]
%     for i in range(1, len(tableau)):
%         valeur = tableau[i]
%         j = i - 1
%         while j >= 0 and tableau[j] > valeur:
%             tableau[j + 1] = tableau[j]
%             j -= 1
%         tableau[j + 1] = valeur
%     return tableau
% depart = [4, 2, 7, 1]
% assert tri_insertion(depart) == [1, 2, 4, 7]
% END-VERIFY
\begin{exercice}{1NSI-AGT-EX-003}{2}{15}
On applique le tri par insertion au tableau \lstinline{[4, 2, 7, 1]}.
\begin{enumerate}
  \item Dérouler l'algorithme : donner le contenu du tableau après chaque itération de la boucle \lstinline{for} externe ($i=1$, $i=2$, $i=3$).
  \item Pour chaque étape, vérifier que le sous-tableau \lstinline{tableau[0:i+1]} est bien trié (c'est l'invariant de boucle du cours).
  \item Le tableau final est-il correctement trié ? Vérifier avec le code.
\end{enumerate}
\end{exercice}
